% Chapter 1. Source pages PDF 17--48 (printed 9--40).
\section*{Mathematical preliminaries from the Introduction}
\begin{lemma}[Initial characterization of the integers]
\label{lem:lbc-intro-integers}
\source{leinster-basic-category-theory:page-0010}
\dcref{0.3}
If a ring $A$ admits exactly one homomorphism $A\to R$ for every ring $R$, then
$A\cong\mathbb Z$.
\end{lemma}
\begin{proof}
The universal property supplies maps $A\to\mathbb Z$ and $\mathbb Z\to A$;
uniqueness forces both composites to be identity maps.
\end{proof}

\begin{lemma}[Tensor-product universal property]
\label{lem:lbc-intro-tensor}
\source{leinster-basic-category-theory:page-0013}
\dcref{0.7}
For vector spaces $U,V$, there is a vector space $U\otimes V$ and bilinear map
$b:U\times V\to U\otimes V$ such that every bilinear map $U\times V\to T$
factors uniquely through a linear map $U\otimes V\to T$.
\end{lemma}
\begin{proof}
Take the free vector space on $U\times V$ and quotient by the subspace
generated by additivity and scalar-linearity relations. The quotient map is
bilinear, and the relations make the induced linear map exist uniquely.
\end{proof}
\section{Categories}
\begin{definition}[Category]
\label{def:lbc-category}
\source{leinster-basic-category-theory:page-0018}
\dcref{1.1.1}
A category $\mathcal A$ consists of a class of objects, a class
$\mathcal A(A,B)$ of maps for every pair of objects, identity maps
$1_A\in\mathcal A(A,A)$, and a composition operation
$\mathcal A(B,C)\times\mathcal A(A,B)\to\mathcal A(A,C)$, $(g,f)\mapsto gf$,
such that $h(gf)=(hg)f$ and $1_Bf=f=f1_A$ whenever the composites are defined.
\end{definition}

\begin{remark}[Notation and composition]
\label{rem:lbc-composition}
\source{leinster-basic-category-theory:page-0024}
\dcref{1.1.10}
Composition is written right-to-left; the same letter may denote an object and
its identity map. A category is small when its objects and maps form sets.
\uses{def:lbc-category}
\end{remark}

\begin{definition}[Isomorphism]
\label{def:lbc-isomorphism}
\source{leinster-basic-category-theory:page-0020}
\dcref{1.1.4}
A map $f:A\to B$ is an isomorphism if there is a map $g:B\to A$ with
$gf=1_A$ and $fg=1_B$.
\uses{def:lbc-category}
\end{definition}

\begin{remark}[Category notation and composites]
\label{rem:lbc-category-notation}
\source{leinster-basic-category-theory:page-0018}
\dcref{1.1.2}
We write $A\in\mathcal A$ for an object, $f:A\to B$ for a map, and $gf$ for
composition. Associativity and identities imply that every composable string of
maps has one unambiguous composite (including the empty composite $1_A$), and
that a diagram commutes when all parallel paths have the same composite.
\uses{def:lbc-category}
\end{remark}

\begin{definition}[Construction: Opposite category]
\label{def:lbc-op-category}
\source{leinster-basic-category-theory:page-0024}
\dcref{1.1.9}
The opposite category $\mathcal A^{\op}$ has the same objects and
$\mathcal A^{\op}(A,B)=\mathcal A(B,A)$, with reversed composition.
\uses{def:lbc-category}
\end{definition}

\begin{definition}[Construction: Product category]
\label{def:lbc-product-category}
\source{leinster-basic-category-theory:page-0024}
\dcref{1.1.11}
The product $\mathcal A\times\mathcal B$ has object pairs and componentwise
maps, identities, and composition.
\uses{def:lbc-category}
\end{definition}

\section{Functors}
\begin{definition}[Functor]
\label{def:lbc-functor}
\source{leinster-basic-category-theory:page-0025}
\dcref{1.2.1}
For categories $\mathcal A,\mathcal B$, a functor $F:\mathcal A\to\mathcal B$
assigns an object $F(A)$ to each $A$ and a map $F(f):F(A)\to F(A')$ to each
$f:A\to A'$, preserving identities and composition:
$F(1_A)=1_{F(A)}$ and $F(gf)=F(g)F(f)$.
\end{definition}

\begin{remark}[Functorial composites]
\label{rem:lbc-functor-composites}
\source{leinster-basic-category-theory:page-0025}
\dcref{1.2.2}
Functoriality makes every composable string of maps have a unique image
composite; functors compose, and each category has an identity functor.
\uses{def:lbc-functor}
\end{remark}

\begin{definition}[Contravariant functor]
\label{def:lbc-contravariant}
\source{leinster-basic-category-theory:page-0030}
\dcref{1.2.10}
A contravariant functor from $\mathcal A$ to $\mathcal B$ is a functor
$\mathcal A^{\op}\to\mathcal B$.
\uses{def:lbc-op-category,def:lbc-functor}
\end{definition}

\begin{definition}[Presheaf]
\label{def:lbc-presheaf}
\source{leinster-basic-category-theory:page-0032}
\dcref{1.2.15}
A presheaf on $\mathcal A$ is a functor $\mathcal A^{\op}\to\Set$.
\uses{def:lbc-contravariant}
\end{definition}

\begin{definition}[Full and faithful functor]
\label{def:lbc-full-faithful}
\source{leinster-basic-category-theory:page-0033}
\dcref{1.2.16}
A functor $F:\mathcal A\to\mathcal B$ is faithful (respectively full) when
each map $\mathcal A(A,A')\to\mathcal B(F(A),F(A'))$ is injective
(respectively surjective).
\uses{def:lbc-functor}
\end{definition}

\begin{remark}[Naturality determines diagonals]
\label{rem:lbc-natural-diagonal}
\source{leinster-basic-category-theory:page-0036}
\dcref{1.3.2}
For a map $f:A\to A'$ the naturality square determines the induced diagonal
$F(A)\to G(A')$; when $f=1_A$ this is the component $\alpha_A$.
\uses{def:lbc-natural-transformation}
\end{remark}

\begin{warning}[Faithfulness warning]
\label{warn:lbc-faithfulness}
\source{leinster-basic-category-theory:page-0033}
\dcref{1.2.17}
Faithfulness concerns each hom-set map for fixed source and target; it does not
assert that distinct maps with different source or target remain distinct.
\uses{def:lbc-full-faithful}
\end{warning}

\begin{definition}[Subcategory]
\label{def:lbc-subcategory}
\source{leinster-basic-category-theory:page-0033}
\dcref{1.2.18}
A subcategory $\mathcal S$ of $\mathcal A$ consists of subclasses of objects
and maps containing identities and closed under composition. It is full when
$\mathcal S(S,S')=\mathcal A(S,S')$ for all its objects.
\uses{def:lbc-category}
\end{definition}

\begin{warning}[Image of a functor]
\label{warn:lbc-functor-image}
\source{leinster-basic-category-theory:page-0033}
\dcref{1.2.19}
The image of a functor need not be closed under composition and therefore need
not itself be a subcategory.
\uses{def:lbc-functor,def:lbc-subcategory}
\end{warning}

\section{Natural transformations}
\begin{definition}[Natural transformation]
\label{def:lbc-natural-transformation}
\source{leinster-basic-category-theory:page-0036}
\dcref{1.3.1}
For functors $F,G:\mathcal A\to\mathcal B$, a natural transformation
$\alpha:F\Rightarrow G$ is a family of maps $\alpha_A:F(A)\to G(A)$ such that
for every $f:A\to A'$,
\[G(f)\,\alpha_A=\alpha_{A'}\,F(f).\]
\uses{def:lbc-functor}
\end{definition}

\begin{definition}[Construction: Functor category]
\label{con:lbc-functor-category}
\source{leinster-basic-category-theory:page-0038}
\dcref{1.3.6}
For categories $\mathcal A,\mathcal B$, the category $[\mathcal A,\mathcal B]$
has functors as objects and natural transformations as maps; identities and
composition are componentwise.
\uses{def:lbc-natural-transformation}
\end{definition}

\begin{definition}[Natural isomorphism]
\label{def:lbc-natural-iso}
\source{leinster-basic-category-theory:page-0039}
\dcref{1.3.10}
A natural transformation is a natural isomorphism when each component is an
isomorphism.
\uses{def:lbc-natural-transformation,def:lbc-isomorphism}
\end{definition}

\begin{definition}[Naturally isomorphic functors]
\label{def:lbc-natural-in-object}
\source{leinster-basic-category-theory:page-0040}
\dcref{1.3.12}
For functors $F,G:\mathcal A\to\mathcal B$, the notation
$F(A)\cong G(A)$ naturally in $A$ means that $F$ and $G$ are naturally
isomorphic.
\uses{def:lbc-natural-iso}
\end{definition}

\begin{remark}[Horizontal composition]
\label{rem:lbc-horizontal-composition}
\source{leinster-basic-category-theory:page-0044}
\dcref{1.3.24}
Natural transformations compose vertically componentwise. Given
$\alpha:F\Rightarrow G$ and $\alpha':F'\Rightarrow G'$, their horizontal
composite $\alpha'\ast\alpha:F'F\Rightarrow G'G$ has component
$(\alpha'\ast\alpha)_A=\alpha'_{G(A)}F'(\alpha_A)=G'(\alpha_A)\alpha'_{F(A)}$;
the two expressions agree by naturality.
\uses{def:lbc-natural-transformation,con:lbc-functor-category}
\end{remark}

\begin{warning}[Isomorphism versus equivalence]
\label{warn:lbc-iso-equivalence}
\source{leinster-basic-category-theory:page-0042}
\dcref{1.3.16}
Object isomorphism and equivalence of categories are distinct uses of the
symbol $\cong$; an equivalence need not be an isomorphism of categories.
\uses{def:lbc-equivalence,def:lbc-isomorphism}
\end{warning}

\begin{lemma}[Pointwise inverse]
\label{lem:lbc-pointwise-inverse}
\source{leinster-basic-category-theory:page-0039}
\dcref{1.3.11}
If $\alpha:F\Rightarrow G$ has all components invertible, then the family
$(\alpha_A)^{-1}$ is a natural transformation $G\Rightarrow F$.
\uses{def:lbc-natural-transformation,def:lbc-isomorphism}
\notready
\end{lemma}
\begin{proof}
The source defers this proof to Exercise 1.3.26.
\end{proof}

\begin{definition}[Equivalence of categories]
\label{def:lbc-equivalence}
\source{leinster-basic-category-theory:page-0042}
\dcref{1.3.15}
An equivalence consists of functors $F:\mathcal A\to\mathcal B$ and
$G:\mathcal B\to\mathcal A$ together with natural isomorphisms
$1_{\mathcal A}\cong GF$ and $FG\cong1_{\mathcal B}$.
\uses{con:lbc-functor-category,def:lbc-natural-iso}
\end{definition}

\begin{definition}[Essentially surjective]
\label{def:lbc-essentially-surjective}
\source{leinster-basic-category-theory:page-0042}
\dcref{1.3.17}
A functor $F:\mathcal A\to\mathcal B$ is essentially surjective on objects if
every $B\in\mathcal B$ is isomorphic to $F(A)$ for some $A\in\mathcal A$.
\uses{def:lbc-isomorphism,def:lbc-functor}
\end{definition}

\begin{proposition}[Characterization of equivalences]
\label{prop:lbc-equivalence-criterion}
\source{leinster-basic-category-theory:page-0042}
\dcref{1.3.18}
A functor is an equivalence if and only if it is full, faithful, and essentially
surjective on objects.
\uses{def:lbc-equivalence,def:lbc-full-faithful,def:lbc-essentially-surjective}
\notready
\end{proposition}
\begin{proof}
The source defers this proof to Exercise 1.3.32.
\end{proof}

\begin{corollary}[Full faithful functors are embeddings]
\label{cor:lbc-full-faithful-reflects-iso}
\source{leinster-basic-category-theory:page-0043}
\dcref{1.3.19}
If $F:\mathcal C\to\mathcal D$ is full and faithful, then $\mathcal C$ is
equivalent to the full subcategory of $\mathcal D$ whose objects are of the
form $F(C)$ for some $C\in\mathcal C$.
\uses{def:lbc-full-faithful,prop:lbc-equivalence-criterion}
\end{corollary}
\begin{proof}
The functor $F$ regarded as a functor into this full subcategory is full,
faithful, and essentially surjective by construction. Apply
Proposition~\ref{prop:lbc-equivalence-criterion}.
\end{proof}
